\section{Layered binary codes and joint support exchanges}\label{sec:codes}

The Leech lift gains points by changing how its finite blocks are used.
The next construction instead keeps the layer geometry fixed and
improves its binary ingredients. We first give the common geometric
argument, then specify the support families and the dense sign codes.
This order separates the dimension-independent proof from the finite
objects that determine the cardinalities.

\subsection{The three-layer construction}
Let $n\ge m\ge32$. Choose a binary code $T\subset\F_2^m$ with minimum
distance at least $\lceil m/4\rceil$, and a support family
$\mathcal B\subseteq\binom{[n]}8$ such that
\begin{equation}\label{eq:support}
 |B\cap B'|\le4\qquad(B\ne B',\ B,B'\in\mathcal B).
\end{equation}
For the standard basis $e_0,\ldots,e_{n-1}$ of $\R^n$, define
\begin{align}
 X_T&=\left\{\frac1{\sqrt m}\sum_{i=0}^{m-1}(-1)^{c_i}e_i:c\in T\right\},
 \label{eq:sign-layer}\\
 X_M&=\left\{\frac1{\sqrt8}\sum_{i\in B}\eps_i e_i:
       B\in\mathcal B,\ \eps_i\in\{\pm1\},\ \prod_{i\in B}\eps_i=1\right\},
 \label{eq:middle-layer}\\
 X_R&=\left\{\frac{\eps e_i+\eps'e_j}{\sqrt2}:
       0\le i<j<n,\ \eps,\eps'\in\{\pm1\}\right\}.
 \label{eq:root-layer}
\end{align}
The final layer consists of the normalized $D_n$ roots.

\begin{proposition}[Three-layer specialization of ERS~\cite{ers}]
\label{prop:ers}
The union $X_T\cup X_M\cup X_R$ is a kissing configuration with
\begin{equation}\label{eq:ers-count}
 |X|=|T|+128|\mathcal B|+2n(n-1).
\end{equation}
\end{proposition}
\begin{proof}
All points are unit vectors. Their support sizes $m$, eight, and two
distinguish the layers. Within each layer, the signs and support recover
the parameters uniquely. On an eight-element support there are $2^7$
even sign patterns; on a two-element support there are four sign
patterns. This proves the count.

In the dense layer, two codewords at distance $d$ have inner product
$1-2d/m\le1/2$. On one middle support, two distinct even sign patterns
differ in at least two positions, giving at most $(8-4)/8=1/2$.
On different middle supports,~\eqref{eq:support} gives at most four
contributions of $1/8$. Two roots with different supports share at most
one coordinate, giving at most $1/2$; on the same support distinct
roots have inner product zero or $-1$.

For the three cross-layer cases, the numbers of overlapping coordinates
give respectively
\begin{equation}\label{eq:ers-cross}
 \ip{X_T}{X_M}\le\sqrt{8/m},\qquad
 \ip{X_T}{X_R}\le\sqrt{2/m},\qquad
 \ip{X_M}{X_R}\le\frac{2}{\sqrt8\sqrt2}=\frac12.
\end{equation}
Here the notation denotes the maximum over the indicated two layers.
The first two bounds are at most $1/2$ because $m\ge32$.
If a middle support extends beyond the first $m$ coordinates, the
overlap is smaller. Thus every pair satisfies~\eqref{eq:kissing}.
\end{proof}

The proof imposes no coupling between particular support words and
particular sign words. Any inputs with the stated parameters can be
combined. In particular, one additional support contributes $128$
points, and an increase in the sign-code size contributes its full
cardinality increment.

\subsection{The support families}
For eight-element supports the identity
\[
 d_H(B,B')=16-2|B\cap B'|
\]
identifies~\eqref{eq:support} with minimum distance eight. The following
finite objects supply the middle layers.

\begin{theorem}\label{thm:supports}
The following constant-weight-code lower bounds hold:
\[
\begin{array}{c|rrrrr}
 n&33&34&37&38&39\\\hline
 A(n,8,8)\ge&1803&1960&2843&3077&3383.
\end{array}
\]
\end{theorem}
\begin{proof}[Finite verification]
The complete binary-mask lists $\mathcal B_n$ are specified in
Table~\ref{tab:certificates}.
For the list of length $n$, every mask is nonnegative and less
than $2^n$, has exactly eight set bits, and occurs once. Every unordered
pair has an intersection of size at most four. Direct pairwise checks
give the distributions in Table~\ref{tab:support-histograms}.
An independent equivalent check enumerates all five-element subsets of
each support: no such subset occurs in two supports. This equivalence
holds because an intersection has size at least five exactly when it
contains a five-subset. In total $56|\mathcal B_n|$ distinct keys are
obtained. These finite checks establish the stated codes.
\end{proof}

\begin{table}[ht]
\centering\small
\caption{Support cardinalities and numbers $h_j$ of unordered pairs
with intersection size $j$. All other intersection counts are zero.}
\label{tab:support-histograms}
\begin{tabular}{rrrrrrr}
\toprule
$n$ & $|\mathcal B_n|$ & $h_0$ & $h_1$ & $h_2$ & $h_3$ & $h_4$\\
\midrule
33 & 1803 & 134282 & 389920 & 680407 & 234225 & 185669\\
34 & 1960 & 168266 & 488949 & 783677 & 281186 & 197742\\
37 & 2843 & 443987 & 1255368 & 1461343 & 621750 & 257455\\
38 & 3077 & 555513 & 1521825 & 1676186 & 706467 & 272435\\
39 & 3383 & 718023 & 1896522 & 1985532 & 820689 & 299887\\
\bottomrule
\end{tabular}

\end{table}

The final support families arise from exchanges in the constant-weight
space. For an incumbent family $\mathcal B$ and a new support $c$, let
$F(c)=\{b\in\mathcal B:|b\cap c|>4\}$. If
$Y\subseteq\binom{[n]}8\setminus\mathcal B$ is internally compatible,
then
\begin{equation}\label{eq:exchange}
 \mathcal B'=(\mathcal B\setminus F(Y))\cup Y,
 \qquad F(Y)=\bigcup_{c\in Y}F(c),
 \qquad |\mathcal B'|-|\mathcal B|=|Y|-|F(Y)|.
\end{equation}
The old--old and new--new pairs are compatible by assumption; all
incompatible old--new pairs have been removed. The union in $F(Y)$
is essential: several inserted words may share deletion costs, allowing
a positive joint exchange even when no single insertion improves the
cardinality.

The final transitions are
\[
\begin{array}{c|rrrrr}
n&33&34&37&38&39\\\hline
\text{parent size}&1802&1959&2840&3076&3382\\
|F(Y)|&9&6&37&6&14\\
|Y|&10&7&40&7&15\\
\text{final size}&1803&1960&2843&3077&3383.
\end{array}
\]
The certificate includes each complete parent and its inserted and
removed sets.
The checker recomputes each conflict union $F(Y)$ and reconstructs
the final family by~\eqref{eq:exchange}. The theorem itself uses the final complete
witnesses, so its verification is independent of choices made during
the search.

\subsection{Dense layers from coset unions}
Let $H\subset\F_2^m$ be linear and let $r_0,\ldots,r_{q-1}$ be coset
representatives. Define
\begin{equation}\label{eq:coset-union}
 T=\bigcup_{a=0}^{q-1}(r_a+H),\qquad
 \delta_{ab}=\min_{h\in H}\wt(r_a+r_b+h)\quad(a\ne b).
\end{equation}
Within one coset, distances between distinct words are the nonzero
weights of $H$. Between two cosets, the difference set is exactly
$r_a+r_b+H$. Thus the $\delta_{ab}$ determine all cross-coset minimum
distances. Positive $\delta_{ab}$ also prove that the cosets are
disjoint, giving $|T|=q2^{\dim H}$.

\begin{proposition}\label{prop:signs}
The generators and representatives in Appendix~\ref{app:codes} define
the sign codes in Table~\ref{tab:signs}.
\end{proposition}
\begin{table}[ht]
\centering
\caption{Parameters of the dense sign codes.}
\label{tab:signs}
\begin{tabular}{rrrrrrl}
\toprule
$m$&$\dim H$&$d(H)$&$q$&$d(T)$&$|T|$&Used in dimension\\
\midrule
32&17&8&1&8&131072&32--37\\
38&14&12&11&10&180224&38\\
39&15&12&10&10&327680&39\\
\bottomrule
\end{tabular}
\end{table}
\begin{proof}[Finite verification]
Binary elimination gives the three displayed ranks. Enumeration of
the spans gives the weight distributions in Appendix~\ref{app:codes},
whose least positive weights are eight, twelve, and twelve. The
length-$32$ code uses the sole representative zero. For lengths $38$
and $39$, enumeration of~\eqref{eq:coset-union} gives
$\delta_{ab}=10$ for all $55$ and $45$ distinct representative pairs,
respectively. The preceding difference-set argument proves the sizes
and minimum distances. Finally, each distance is at least
$\lceil m/4\rceil$, as required by Proposition~\ref{prop:ers}.
\end{proof}

The length-$32$ code is the Cheng--Sloane
$[32,17,8]$ code~\cite{chengsloane}, in the generator convention of
Grassl's table~\cite{grassl32}. The other two kernels have a common
quadratic-residue origin. Let
\begin{align}
 g(x)={}&x^{23}+x^{19}+x^{18}+x^{14}+x^{13}+x^{12}+x^{10}+x^9
 \nonumber\\
 &\quad+x^7+x^6+x^5+x^3+x^2+x+1\in\F_2[x].\label{eq:qr}
\end{align}
Take the coefficient vectors of $x^ig(x)$, $0\le i<24$, as
length-$47$ binary words, and append their parity at coordinate $47$.
This specifies the extended quadratic-residue code in the chosen
convention~\cite[Chapter~16]{macwilliams}. Shorten at coordinates
$40,\ldots,47$, retaining coordinates $0,\ldots,39$ in order, to
obtain a dimension-$16$ code $H_{40}$. Shortening $H_{40}$ at coordinate
zero gives the printed $H_{39}$; shortening at coordinates zero and
one gives the printed $H_{38}$. Here shortening imposes zeros before
deleting coordinates. Exact row-space comparisons verify these
identities and, in the displayed retained-coordinate conventions,
\begin{equation}\label{eq:nesting}
 H_{38}=\{(h_0,\ldots,h_{37}):(h_0,\ldots,h_{37},0)\in H_{39}\}.
\end{equation}
These reconstructions give both kernels a precise common source.
The distance proof uses their explicit generators and the affine
representatives rather than an assumed distance of a named mother code.

\begin{proof}[The five layered-code constructions]
For each $n\in\{33,34,37,38,39\}$, combine the support family of
Theorem~\ref{thm:supports} with the sign code of
Proposition~\ref{prop:signs}. Proposition~\ref{prop:ers} gives
Table~\ref{tab:layers}. The dimension-37 and dimension-38 counts are
improved further in Sections~\ref{sec:signed} and~\ref{sec:completion}.
\end{proof}

\begin{table}[ht]
\centering
\caption{The three contributions to the five code-based configurations.}
\label{tab:layers}
\begin{tabular}{rrrrr}
\toprule
$n$&$|X_T|$&$|X_M|$&$|X_R|$&Total\\
\midrule
33&131072&230784&2112&363968\\
34&131072&250880&2244&384196\\
37&131072&363904&2664&497640\\
38&180224&393856&2812&576892\\
39&327680&433024&2964&763668\\
\bottomrule
\end{tabular}
\end{table}

At the first three lengths, the increase relative to the cited
comparison is $128$ times the increase in support size, respectively
$15$, $24$, and $11$. At the last two lengths, both the larger support
families and the explicit coset unions enter the final count.
The same geometric argument accommodates all five inputs, while
their different dense-layer distances lead to the different local
geometries considered next.

\subsection{A new length-32 support family}
The same construction applies to a verified family of 1676 eight-subsets
of $[32]$. There are $1,403,650$ unordered support pairs. Their intersection
histogram, in orders zero through four, is
\[
(112077,318935,601033,196122,175483).
\]
No larger intersection occurs. Together with the same $[32,17,8]$ dense
code, Proposition~\ref{prop:ers} gives
\[
K(32)\ge 2^{17}+128(1676)+2(32)(31)=347584.
\]
The family was obtained from Echols's public 1667-word seed~\cite{echols} through joint
exchanges. The final witness, rather than successful replay of a randomized
search, is the finite input to this proof. The 2843-word length-37 family
above is also improved to 2845 words. Section~\ref{sec:signed} uses this
larger family and changes the sign patterns, giving a further gain.
